\section*{Laplace Transformation Anwendungen Aufgaben mit L�sungen}\label{MA2-laplacetransformation-und-differentialgleichungen-loesungen}

\ifintegriereaufgaben Erkl�rungen und Beispiele, die zum Verst�ndnis der L�sungswege hilfreich sind, finden Sie im Skript\wennreferenzvorhanden{MA2-laplacetransformation-und-differentialgleichungen}{ (Abschnitt \abschnittsverweis{MA2-laplacetransformation-und-differentialgleichungen})}. \fi

\begin{exercisebox}[Laplace-Transformation Rechenregeln] % Papula 2, VI, Abschnitt 2, Aufgabe 1, 2, 3, 5, 6, 7, 9, 10
	Die Aufgaben finden Sie in Papula 2.VI.2 ab Seite \(677\) und die L�sungen ab Seite \(767\).
	\begin{itemize}
		\item[(a)] Papula, Aufg. 1a) $\quad \mathcal L_f(s) = \frac{24}{s^4} - \frac{2}{s^3} + \frac{2}{s^2} $
	\item[(b)] Papula, Aufg. 1.b) $\quad \mathcal L_f(s) = C\left(\frac1s - \frac{1}{s+\lambda}\right) $
\item[(c)] Papula, Aufg. 2.a) $\quad \mathcal L_f(s) = \frac{29160}{s^6} $
		\item[(d)] Papula, Aufg. 2.b) $\quad \mathcal L_f(s) = \frac{s}{s^2+16}$
	\item[(e)] Papula, Aufg. 2.d) $\quad \mathcal L_f(s) = \frac{s^2+2\omega^2}{s(s^2+4\omega^2)} $
\item[(f)] Papula, Aufg. 3.b) $\quad \mathcal L_f(s) = \frac{2\cdot e^{-4s}}{s^3} $
		\item[(g)] Papula, Aufg. 3.c) $\quad \mathcal L_f(s) = \frac{e^{-bs}}{s-1} $
	\item[(h)] Papula, Aufg. 3.a) $\quad \mathcal L_f(s) = \frac{s}{s^2+1} $ (Achtung: Verschiebungsformel nicht anwendbar, weil keine Verschiebung nach rechts vorliegt!)
\item[(i)] Papula, Aufg. 5.a) $\quad \mathcal L_f(s) = \frac{s-3}{(s-3)^2+4} $
\item[(j)] Papula, Aufg. 5.b) $\quad \mathcal L_f(s) = \frac{A\omega}{(s+\delta)^2+\omega^2} $
		\item[(k)] Papula, Aufg. 5.c) $\quad \mathcal L_f(s) = \frac{1}{s-3\ln(2)}$
		\end{itemize}
\end{exercisebox}


\begin{exercisebox}[Laplace-Transformation und Differentialgleichungen]
	Betrachten Sie die Differentialgleichung $ y'' + \omega^2 y = 0,\; \omega>0$ mit den beiden Anfangswerten
	\begin{enumerate}
		\item $y(0) = 1 \land y'(0) = 0$
		\item $y(0) = 0 \land y'(0) = 1$
	\end{enumerate}
	\begin{itemize}
		\item[(a)] Berechnen Sie die L�sung der beiden Anfangswertprobleme.
		\item[(b)] Berechnen Sie die Laplace-Transformierte $Y(s)$ der Differentialgleichung \( y'' + \omega^2\, y = 0\). Nutzen Sie dabei die Ableitungsregeln.
		\item[(c)] Bestimmen Sie, ausgehend von den speziellen L�sungen der Anfangswertprobleme im Bildbereich, die L�sungen im Zeitbereich. Verwenden Sie f�r die R�cktransformationen eine Tabelle.
	\end{itemize}
\vspace{0.3cm}
\textbf{L�sung:} \newline
	\begin{itemize}
		\item[(a)] Die L�sungen lauten
			\[y_1(t) = \cos(\omega t), \; y_2(t) = 1/\omega \cdot \sin(\omega t)\]
		\item[(b)] Laplace-Transformation
			\begin{eqnarray*}
				\mathcal L\left(y'' + \omega^2 y\right)(s) &=& \mathcal L\left(y''\right) (s) + \omega^2\mathcal L\left( y\right)(s) \\
				&=& s^2\mathcal L\left( y\right) (s) - s y(0) + \omega^2 \mathcal L\left( y\right) (s) \\
				&=& s^2 Y (s) - s y(0) - y'(0) + \omega^2 Y(s) \stackrel{!}{=} 0\,.
			\end{eqnarray*}
		\item[(c)] Erstes Anfangswertproblem
			\[ s^2 Y_1(s) - s+ \omega^2 Y_1(s) = 0 \Leftrightarrow Y_1(s) = \frac{s}{s^2+\omega^2} \Rightarrow y_1(t) = \cos(\omega t)\,.\]
		Zweites Anfangswertproblem
		\[ s^2 Y_2(s) - 1 + \omega^2 Y_2(s) = 0 \Leftrightarrow Y_2(s) = \frac{1}{s^2+\omega^2} \Rightarrow y_2(t) = \frac{1}{\omega}\cdot \sin(\omega t)\,.\]
\end{itemize}
\end{exercisebox}

\begin{exercisebox}[Laplace-Transformation und Differentialgleichungen]
	L�sen Sie die folgenden Systeme mit Hilfe der Laplace-Transformation
	\begin{eqnarray*}
		(a)\; \begin{cases}
			x(0) = y(0) =  0, \quad & t= 0\,,\\
			\dot x + 2 y = e^t, \quad & t> 0\,,\\
			\dot y + 2 x = e^{-t}, \quad & t> 0\,.
		\end{cases} \quad\quad
		(b)\; \begin{cases}
			x(0) = y(0) =  0, \quad & t= 0\,,\\
			\dot x(0) = \dot y(0) =  0, \quad & t= 0\,,\\
			\ddot x + \dot y + 3 y = 1, \quad & t> 0\,,\\
			\ddot y -4 \dot x + 3y = 0, \quad & t> 0\,.
		\end{cases} \quad\quad
	\end{eqnarray*}
\vspace{0.3cm}
\textbf{L�sung:} \newline
	\begin{itemize}
		\item[(a)]
			Anwendung der Laplace-Transformation und Einsetzen der Anfangsbedingungen ergibt:
			\[
		\begin{cases}
			s X(s) + 2 Y(s) = \displaystyle{\frac{1}{s-1}}\,,\\
			2 X(s) + sY(s) = \displaystyle{\frac{1}{s+1}}\,.
	\end{cases} \]

%			Und nach Einsetzen der homogenen Anfangsbedingungen erhalten wir das folgende lineare Gleichungssystem:
%			\[
%				\left( \begin{array}{cc} s & -2 \\ -2 & s\end{array}\right) \left( \begin{array}{c} X \\Y \end{array}\right)(s) = \frac{1}{s^2-1}\left( \begin{array}{c} s+1 \\ s-1 \end{array}\right) \,.
%\]
	Aufl�sen\footnote{ Entweder durch Einsetzen der Gleichungen ineinander oder durch Anwendung der inversen Matrix!} und Partialbruchzerlegung liefert
			\[
		\begin{cases}
			X(s) = \displaystyle{\frac{s^2 - s + 2}{(s^2-1)(s^2-4)}} = \displaystyle{\frac{1}{3}}\left( \displaystyle{\frac{2}{s+1}} - \displaystyle{\frac{1}{s-1}} - \displaystyle{\frac{2}{s+2}} + \displaystyle{\frac{1}{s-2}}\right)\,,\\
			Y(s) = \displaystyle{\frac{s^2 - 3s - 2}{(s^2-1)(s^2-4)}} = \displaystyle{\frac{1}{3}}\left( \displaystyle{\frac{1}{s+1}} + \displaystyle{\frac{2}{s-1}} - \displaystyle{\frac{2}{s+2}} - \displaystyle{\frac{1}{s-2}}\right)\,.
	\end{cases} \]
	R�cktransformation liefert die L�sung des Systems:
			\[
		\begin{cases}
			x(t) =  \displaystyle{\frac{1}{3}}\left( 2 e^{-t} - e^{t} - 2 e^{-2t} + e^{2t} \right)\,,\\
			y(t) =  \displaystyle{\frac{1}{3}}\left( e^{-t} +2 e^{t} - 2 e^{-2t} - e^{2t} \right)\,.
	\end{cases} \]

		\item[(b)]
			Laplace-Transformation des Systems unter Bere�cksichtigung der Anfangsbedingungen:
			\[
		\begin{cases}
			\left( s^2 + 3\right) X(s) + s Y(s) = \displaystyle{\frac{1}{s}}\,,\\
			-4s X(s) + \left( s^2 + 3\right) Y(s) = 0\,.
	\end{cases} \]
	Aufl�sen und Partialbruchzerlegung liefert
			\[
		\begin{cases}
		X(s) = \displaystyle{\frac{s^2 +3}{s(s^2+9)(s^2+1)}} =\displaystyle{\frac{1}{3 s}} - \displaystyle{\frac{1}{12(s^2+9)}} - \displaystyle{\frac{s}{4(s^2+1)}}\,,\\
	Y(s) = \displaystyle{\frac{4}{(s^2+9)(s^2+1)}} = -\displaystyle{\frac{1}{2(s^2+9)}} + \displaystyle{\frac{1}{2(s^2+1)}} \,.
	\end{cases} \]
	R�cktransformation liefert die L�sung des Systems:
			\[
		\begin{cases}
			x(t) =  \displaystyle{\frac{1}{3}}- \displaystyle{\frac{1}{12}}\cos(3t)-\frac14\cos(t)\,,\\
			y(t) =  -\displaystyle{\frac{1}{6}}\sin(3t) + \displaystyle{\frac12} \sin(t)\,.
	\end{cases} \]
	\end{itemize}
\end{exercisebox}

